\section{Results: Audit Findings and Paraformer Experiment}
The audit produces three results. First, each of the thirteen
constant-coefficient assumptions has a documented failure regime in ordinary
engineering literature. Second, the failure regimes are not isolated anomalies:
they form a repeatable pattern in which a coefficient becomes a function of
state variables that the basic pedagogical model has suppressed. Third, when a
coefficient is written as a function rather than a number, the missing term is
not arbitrary. It is the derivative or memory contribution that follows from
the product rule, from mutual-coupling kinematics, or from the more general
causal response functional.

The forthcoming empirical component of the paper will test the paraformer
case, because the paraformer sits at the point where the inductive,
permeability and saturation entries of the audit become an experimentally
measurable transformer architecture.

%The present version reserves the exact
%parameter and test locations so that later experimental data can be inserted
%without changing the theoretical structure of the article.


\subsection{Metglas AMCC-1000 parameter description}

The paraformer test section is instantiated here as an orthogonal Metglas
AMCC-1000 centre \cite{amcc1000}. The Metglas centre is treated as an instrumented, resonant,
two-axis magnetic coupler rather than as an ordinary closed transformer. The
first U-core, denoted \(U_A\), is aligned with the \(X\)-axis of the device.
The second U-core, \(U_B\), is rotated by \(90^\circ\) and aligned with the
\(Y\)-axis. The coupling coefficient \(k\) is therefore not assumed to be
transformer-like: the default orthogonal model uses \(k=0.08\), and the
transfer sweep separately reports stronger hypothetical values up to
\(k=0.95\).

The analysis is a conventional electromagnetic simulation. It does not assume
any non-classical energy source. Every reported watt, volt-ampere and var comes
from the coupled RLC state equations, phasor extraction and the matched-load
Thevenin bound. The calculations are intended to guide low-power bench
experiments before any high-voltage or high-current hardware is attempted.

The physical dimensions used for the core are taken from the AMCC-1000 product
drawing for Metglas Alloy 2605SA1 \cite{metglas2605sa1}. The material properties below are the
manufacturer's general 2605SA1 data and are used only as engineering reference
values; the dynamic simulator uses the measured/design inductance,
capacitance and damping values in the CSV files.

\begin{table}[H]
\centering
\textbf{Table 3.} \emph{AMCC-1000 dimensions and magnetic area values used in the simulation.}

\smallskip
\begin{tabular}{llr}
\toprule
Symbol & Meaning & Value \\
\midrule
\(a\) & core build / limb build & \SI{33.0}{\milli\metre} \\
\(b\) & window width & \SI{40.0}{\milli\metre} \\
\(c\) & window length & \SI{105.0}{\milli\metre} \\
\(d\) & core height & \SI{85.0}{\milli\metre} \\
\(e\) & core width & \SI{106.0}{\milli\metre} \\
\(f\) & core length & \SI{171.0}{\milli\metre} \\
\(A_e\) & net magnetic cross-sectional area & \SI{23.0}{\centi\metre\squared} \\
\(W_a\) & window area & \SI{42.0}{\centi\metre\squared} \\
\(A_p\) & area product & \SI{966}{\centi\metre\tothe{4}} \\
\(\eta_p\) & pack factor & \SI{82}{\percent} \\
\(m\) & core mass & \SI{7109}{\gram} \\
\bottomrule
\end{tabular}
\end{table}

\begin{table}[H]
\centering
\textbf{Table 4.} \emph{Metglas 2605SA1 material reference properties used for design context.}

\smallskip
\begin{tabular}{llr}
\toprule
Property & Symbol & Value \\
\midrule
Saturation induction, as cast & \(B_s\) & \SI{1.56}{\tesla} \\
Maximum DC permeability, annealed & \(\mu_{\max}\) & \(6.0\times10^5\) \\
Maximum DC permeability, as cast & \(\mu_{\max}\) & \(4.5\times10^4\) \\
Saturation magnetostriction & \(\lambda_s\) & \SI{27}{ppm} \\
Electrical resistivity & \(\rho_e\) & \SI{130}{\micro\ohm\centi\metre} \\
Curie temperature & \(T_C\) & \SI{395}{\celsius} \\
Ribbon thickness & \(t_r\) & \SI{23}{\micro\metre} \\
Density & \(\rho_m\) & \SI{7.18}{\gram\per\centi\metre\cubed} \\
Lamination factor & -- & \(>84\%\) \\
Continuous service temperature & -- & \SI{150}{\celsius} \\
Crystallization temperature & -- & \SI{510}{\celsius} \\
\bottomrule
\end{tabular}
\end{table}

The dynamic model has one modal current for \(U_A\) and one for \(U_B\). In the
current design each core mode is resonated at \SI{128}{\kilo\hertz}, uses an
eight-turn reference winding, and uses a \SI{2}{\milli\metre} equivalent gap.
The measured/design inductance and capacitance values are
\begin{equation}
L_A=L_B=\SI{92.051753}{\micro\henry}, \qquad
C_A=C_B=\SI{16.795319}{\nano\farad}.
\end{equation}
The ideal resonance condition is
\begin{equation}
f_0 = \frac{1}{2\pi\sqrt{LC}}, \qquad
C = \frac{1}{\omega^2 L}.
\end{equation}
At \SI{128}{\kilo\hertz}, the inductive reactance is
\begin{equation}
X_L = \omega L = \SI{74.032413}{\ohm}.
\end{equation}
With the selected loaded core quality factor \(Q_c=12\), the dynamic damping
resistance used by the modal solver is
\begin{equation}
R_d = \frac{X_L}{Q_c} = \SI{6.169368}{\ohm}.
\end{equation}

The coupled state equation solved by the simulator is
\begin{align}
\mathbf{L}\frac{d\mathbf{i}}{dt} &= \mathbf{v}_s(t)-\mathbf{R}\mathbf{i}-\mathbf{v}_C,\\
\frac{d\mathbf{v}_C}{dt} &= \mathbf{C}^{-1}\mathbf{i}.
\end{align}
The mutual inductance between any two modes \(i\) and \(j\) is
\begin{equation}
M_{ij}=k_{ij}\sqrt{L_iL_j}.
\end{equation}
For the orthogonal Metglas pair in the default model,
\begin{equation}
k_{AB}=0.08, \qquad M_{AB}=\SI{7.364140}{\micro\henry}.
\end{equation}

\subsection{Metglas AMCC-1000 dynamic transformer test}

The simulator outputs define four Metglas-specific waveform groups. First, the
raw modal currents \(i_A(t)\) and \(i_B(t)\) are the simulated resonant
currents in the two AMCC windings. Second, \(v_{C,A}(t)\) and \(v_{C,B}(t)\)
are the modal tank-capacitor voltages. Third, the direct inter-core induced
voltages are
\begin{align}
v_{B\leftarrow A}(t) &= M_{AB}\frac{di_A}{dt},\\
v_{A\leftarrow B}(t) &= M_{AB}\frac{di_B}{dt}.
\end{align}
Fourth, the practical monitor signals at the centre interface are
\begin{align}
v_{\mathrm{in},A4\rightarrow U_A}(t)
&= \sum_{j\in A4} M_{A,j}\frac{di_j}{dt},\\
v_{\mathrm{out},U_B\rightarrow B4}(t)
&= \sum_{j\in B4} M_{j,B}\frac{di_B}{dt}.
\end{align}
These two monitor voltages are not ideal transformer terminal voltages. They
are observable model signals that separate the A4-driven input into \(U_A\)
from the \(U_B\)-driven output delivered toward the B4 branch set.

The instantaneous mutual power terms reported in the waveform CSV are
\begin{align}
p_{A\rightarrow B}(t)&=v_{B\leftarrow A}(t)i_B(t),\\
p_{B\rightarrow A}(t)&=v_{A\leftarrow B}(t)i_A(t),\\
p_{\mathrm{net},A\rightarrow B}(t)&=p_{A\rightarrow B}(t)-p_{B\rightarrow A}(t).
\end{align}
Because inductive coupling stores and returns magnetic energy each cycle, the
instantaneous curves are not by themselves a load-power claim. The matched-load
transfer limit is computed separately as
\begin{equation}
V_{\mathrm{oc},2}=\omega M I_{1,\mathrm{rms}},
\qquad
P_{\max}=\frac{V_{\mathrm{oc},2}^{2}}{4R_{d,2}}.
\end{equation}

Flux density is reconstructed from the modal current by
\begin{equation}
B(t)=\frac{L i(t)}{N A_e}.
\end{equation}
For sinusoidal verification the equivalent voltage relation is
\begin{equation}
V_{\mathrm{rms}}=4.44288 f N A_e B_{\mathrm{peak}}.
\end{equation}

The Python run uses 34 coupled modes and writes 129 time samples over the last
eight cycles at \SI{128}{\kilo\hertz}. Its dynamic Metglas summary is:

\begin{table}[H]
\centering
\textbf{Table 5.} \emph{Dynamic Metglas centre signals from the full coupled ladder model.}

\smallskip
\begin{tabular}{lrrl}
\toprule
Metric & RMS or average & Peak & Unit \\
\midrule
\(U_A\) core current & 2.478628 & 3.449532 & A \\
\(U_B\) core current & 3.367880 & 4.758976 & A \\
\(U_A\) tank voltage & 183.020191 & 255.414153 & V \\
\(U_B\) tank voltage & 249.530296 & 352.351148 & V \\
\(U_A\) flux density & 0.012400 & 0.017257 & T \\
\(U_B\) flux density & 0.016849 & 0.023808 & T \\
Direct \(U_A\rightarrow U_B\) mutual voltage & 14.275697 & 19.908115 & V \\
Input \(A4\rightarrow U_A\) mutual voltage & 21.629079 & 30.577635 & V \\
Output \(U_B\rightarrow B4\) mutual voltage & 54.305273 & 76.609354 & V \\
Average \(U_A\rightarrow U_B\) mutual power & 10.267611 & 58.458069 & W \\
Average net \(U_A\rightarrow U_B\) imbalance & 21.344931 & 22.945421 & W \\
\bottomrule
\end{tabular}
\end{table}

\begin{figure}[H]
\centering
\includegraphics[width=0.98\textwidth]{../figures/metglas_core_currents.png}
\caption{Simulated \(U_A\) and \(U_B\) AMCC-1000 modal currents during the last eight \SI{128}{\kilo\hertz} cycles.}
\end{figure}

\begin{figure}[H]
\centering
\includegraphics[width=0.98\textwidth]{../figures/metglas_core_voltages.png}
\caption{Metglas tank voltages and direct \(U_A\rightarrow U_B\) mutual voltage.}
\end{figure}

\begin{figure}[H]
\centering
\includegraphics[width=0.98\textwidth]{../figures/metglas_input_output_signals.png}
\caption{Input/output monitor signals: summed A4 input into \(U_A\), summed \(U_B\) output toward B4, and the direct \(U_A\rightarrow U_B\) induced voltage.}
\end{figure}

\begin{figure}[H]
\centering
\includegraphics[width=0.98\textwidth]{../figures/metglas_transferred_power_time.png}
\caption{Instantaneous mutual power exchange. The oscillatory terms include stored magnetic energy exchange; average load power must be interpreted through the matched-load model.}
\end{figure}

The current full-ladder dynamic state has modest core current because the
Metglas windings are not directly driven in the same way as the separately
optimized near-\SI{10}{\kilo\watt} paraformer model. The matched-load Thevenin
limits are:

\begin{table}[H]
\centering
\textbf{Table 6.} \emph{Metglas matched-load transfer limits for the current \(k=0.08\) model.}

\smallskip
\begin{tabular}{lrrrr}
\toprule
Case & \(I_1\) (A rms) & \(V_{\mathrm{oc},2}\) (V rms) & \(R_{d,2}\) (ohm) & \(P_{\max}\) (W) \\
\midrule
Dynamic \(U_A\rightarrow U_B\) & 2.475577 & 14.661833 & 6.169368 & 8.711158 \\
Dynamic \(U_B\rightarrow U_A\) & 3.369146 & 19.954079 & 6.169368 & 16.134768 \\
Target 50 mT \(U_A\rightarrow U_B\) & 7.067092 & 41.855512 & 6.169368 & 70.991223 \\
Target 50 mT \(U_B\rightarrow U_A\) & 7.067092 & 41.855512 & 6.169368 & 70.991223 \\
\bottomrule
\end{tabular}
\end{table}

The sweep file shows that a conventional, strong-coupling value \(k=0.95\)
would give approximately \SI{10010.871747}{\watt} at the 50 mT target current. The
orthogonal \(k=0.08\) arrangement reaches the same nominal
\SI{10}{\kilo\watt} matched-load target only by raising the design flux density
and changing the winding count. The optimized orthogonal design point is:

\begin{table}[H]
\centering
\textbf{Table 7.} \emph{Separate near-\SI{10}{\kilo\watt} orthogonal paraformer design point.}

\smallskip
\begin{tabular}{lr}
\toprule
Quantity & Value \\
\midrule
Frequency & \SI{128}{\kilo\hertz} \\
Assumed orthogonal coupling & \(k=0.08\) \\
Gap & \SI{2.0}{\milli\metre} \\
Required peak flux density & \SI{0.593428}{\tesla} \\
Turns per winding & 12 \\
Primary winding current & \SI{55.917426}{\ampere} rms \\
Primary winding voltage & \SI{9.314329}{\kilo\volt} rms \\
Secondary open-circuit voltage & \SI{745.146358}{\volt} rms \\
Matched secondary load & \SI{13.881077}{\ohm} \\
Matched secondary current & \SI{26.840365}{\ampere} rms \\
Self inductance & \SI{0.207116}{\milli\henry} \\
Mutual inductance & \SI{16.569316}{\micro\henry} \\
Resonance capacitance & \SI{7.464586}{\nano\farad} \\
Circulating reactive power & \SI{0.520833}{\mega VAR} \\
Parallel AWG10-equivalent Litz bundles & 4 per winding \\
Active cable per winding & \SI{40.8}{\metre} \\
Procurement cable, both windings & \SI{93.84}{\metre} \\
\bottomrule
\end{tabular}
\end{table}

\begin{figure}[H]
\centering
\includegraphics[width=0.98\textwidth]{../figures/metglas_transfer_power_sweep.png}
\caption{Matched-load transfer bound as a function of coupling coefficient \(k\).}
\end{figure}

\begin{figure}[H]
\centering
\includegraphics[width=0.92\textwidth]{../figures/orthogonal_paraformer_10kw_design.png}
\caption{Scaled summary of the separate near-\SI{10}{\kilo\watt} orthogonal design point.}
\end{figure}

The corresponding bench procedure is as follows. First, mount the AMCC-1000
cores in the \(0^\circ/90^\circ\) orthogonal geometry and measure each winding
with an LCR meter and a vector network analyzer. Record \(L_A\), \(L_B\),
winding resistance, self-resonant frequency and the fitted parallel
capacitance. Second, drive \(U_A\) through a current-limited RF source at low
voltage and measure the open-circuit induced voltage on \(U_B\). Estimate
\begin{equation}
k = \frac{V_{\mathrm{oc},B}}{\omega I_A\sqrt{L_A L_B}}.
\end{equation}
Repeat the reciprocal \(U_B\rightarrow U_A\) test. Third, install the
calculated resonant capacitance and trim with a VNA until the phase of the
input impedance crosses zero at \SI{128}{\kilo\hertz}. Measure the loaded
quality factor from
\begin{equation}
Q_L=\frac{f_0}{f_2-f_1}.
\end{equation}
Use the measured \(Q_L\) to update \(R_d=\omega L/Q_L\) in the simulator.

At low power, probe \(i_A(t)\), \(i_B(t)\), \(v_{C,A}(t)\), and
\(v_{C,B}(t)\) with isolated high-voltage probes and RF current probes. Compare
measured RMS and phase with the simulated Metglas signal summary. For rectified
output experiments on the B side, use a resonant matched load or an RF-rated
SiC diode bridge followed by a controlled DC load. Diode junction capacitance
must be included in the measured output tuning because at
\SI{128}{\kilo\hertz} it can shift the B-side resonance and alter \(Q_L\).
Report real DC output power as
\begin{equation}
P_{\mathrm{DC}}=\langle v_{\mathrm{DC}}(t)i_{\mathrm{DC}}(t)\rangle
\end{equation}
and compare it against RF input real power, not against circulating var.

The present dynamic model shows that the orthogonal centre is strongly
reactive. The plotted currents, voltages and mutual-power waveforms are
coherent at \SI{128}{\kilo\hertz}. The near-\SI{10}{\kilo\watt} case is a
separate high-stress extrapolation in which the flux density is increased to
\SI{0.593428}{\tesla}, the winding count is changed to 12 turns, and the
circulating reactive power rises to roughly \SI{0.520833}{\mega VAR}. That point
is below the nominal 2605SA1 saturation induction, but it is high enough that
core loss, thermal rise, insulation stress and layout parasitics become
first-order experimental variables.
